% =============================================================================
% Project : THE_PHILOSOPHY_OF_VIABILITY
% Book    : Book I — Ontology of Viability
% Book ID : VS.B01
%
% File    : 2_MAINMATTER/PART_I_IDENTITY_AND_MANIFESTATION/CHAPTER_02_INVARIANCE/sections/section_11_synthesis_of_invariance.tex
%
% Section : 11 — Synthesis of Invariance
% Chapter : 02 — Invariance
%
% Purpose :
%   To consolidate the ontological account of invariance, define
%   its proper conceptual and inferential scope, distinguish its
%   analytical dimensions, and identify the conclusions and
%   unresolved questions established by the chapter.
%
% Conceptual Scope :
%   Invariance; determination; transformation; preservation;
%   identity; numerical identity; persistence; continuity;
%   stability; immutability; forms and domains of invariance;
%   formal and ontological invariance; positive and negative
%   inference; unresolved ontological questions.
%
% =============================================================================

\section{Synthesis of Invariance}
\label{sec:synthesis-of-invariance}

\subsection{The Position of Invariance Within the Ontology of Identity}

The preceding sections have examined invariance as an
ontological problem concerning the preservation of
determinations under transformation. The analysis has
distinguished invariance from identity and has examined
its relations to determination, difference, structure,
composition, continuity, and the domains within which
preservation may be investigated.

The resulting account establishes invariance as a
significant but limited concept. It identifies a particular
relation between a determination and a transformation:
a specified determination remains preserved under a
specified transformation, within a specified scope.

This relation is significant because transformation does
not necessarily affect every determination in the same
way. Some determinations may remain preserved while
others are altered, replaced, reorganized, or left
undetermined by the analysis. Invariance makes this
differentiation conceptually explicit.

Nevertheless, the preservation of a determination does
not constitute the identity of the existent in which
that determination occurs. Nor does the alteration of
a determination independently establish the loss of
that identity.

Invariance must therefore be situated within the broader
ontology of identity without being made its foundation,
definition, or exhaustive criterion. Its contribution
concerns the analysis of preservation within transformation,
not the independent resolution of every question
concerning the existent.

\subsection{The Definition and Scope of Invariance}

On the basis of the preceding analysis, invariance is
defined as follows:

\begin{quote}
Invariance is the preservation of a specified
determination of an existent under a specified
transformation, relative to a defined domain,
respect, criterion, and range of validity.
\end{quote}

This definition expresses a relational and qualified
concept of preservation. It does not attribute
unrestricted sameness to an existent, nor does it
presume that every determination remains unchanged
whenever transformation occurs.

Several elements are essential to its interpretation.

First, invariance concerns a specified determination.
The determination must be distinguished from the
existent considered as a whole. The preservation of
one determination does not establish the preservation
of all the determinations belonging to that existent.

Second, invariance is asserted under a specified
transformation. A determination may remain preserved
under one transformation and be altered under another.
An invariance claim is therefore not meaningful as
an unrestricted attribution detached from transformation.

Third, the claim is made within a defined domain and
respect. These delimit the field and aspect in which
preservation is being examined. A claim concerning
one domain or respect cannot automatically be extended
to another.

Fourth, preservation must be assessed according to
a criterion. The criterion specifies what counts as
preservation rather than alteration, approximation,
resemblance, or another relation that may be relevant
to comparison.

Finally, the claim has a range of validity. Its
conclusion cannot exceed the transformations,
conditions, and determinations for which preservation
has been established.

These qualifications are not incidental additions
to an otherwise unrestricted concept. They determine
the content and scope of an invariance claim.

\subsection{Preservation and Identity}

The central ontological conclusion of this chapter
is that preservation and identity must remain
analytically distinct.

Identity concerns the existent as such. Invariance
concerns the preservation of a specified determination
under transformation. These concepts address different
questions, even when they are examined in relation
to the same existent.

Distinct existents may possess equivalent or
indistinguishable determinations. Their agreement
with respect to a specified invariant does not
establish their numerical identity.

Conversely, an existent may undergo transformations
in which previously identified invariant determinations
are altered or cease to be preserved. Such alteration
does not independently establish that the existent
has lost its identity.

The relation between identity and invariance must
therefore avoid two symmetrical reductions:

\begin{itemize}
    \item Identity cannot be reduced to the preservation
          of one or more determinations.
    \item The alteration or absence of an identified
          invariant cannot be treated as sufficient
          evidence for the absence or loss of identity.
\end{itemize}

This analytical distinction does not require identity
and invariance to be treated as unrelated in the
concrete existent. They may be examined together,
and findings concerning preservation may contribute
to arguments about identity. But the relevance of
such findings must be independently established.

The preservation of a determination is not identical
to the identity of that which possesses it.

\subsection{Transformation as a Condition of Invariance}

Invariance cannot be understood independently of
transformation. A preservation claim identifies what
remains preserved under a transformation; without
specifying the transformation, the scope of the
claim remains undetermined.

The same determination may be preserved under one
class of transformations and altered under another.
Consequently, invariance is transformation-relative
rather than an unqualified property attributed to
a determination in isolation.

This relativity does not make invariance subjective
or arbitrary. It means that the validity of a claim
must correspond to the transformation and conditions
under which preservation is asserted.

Transformation may also involve preservation and
alteration simultaneously. One determination may
remain preserved while another changes. A structural
relation may be maintained while components are
replaced; a compositional feature may be preserved
while the organization of the whole changes; or
a relational determination may remain preserved
while other aspects of the existent are modified.

Such possibilities must be assessed according to
the determinations and respects involved. Neither
preservation nor alteration should be assumed to
describe the existent exhaustively.

\subsection{The Analytical Requirements of an Invariance Claim}

The analysis of this chapter establishes that an
invariance claim must be sufficiently specified
before its validity and implications can be assessed.

At a minimum, the following elements must be
distinguished:

\begin{enumerate}
    \item \textbf{Subject of the Claim:}
          The existent, configuration, or analytical
          subject whose determination is being examined.

    \item \textbf{Specified Determination:}
          The particular determination alleged to remain
          preserved. It must not be confused with the
          existent considered as a whole.

    \item \textbf{Transformation:}
          The transformation under which preservation
          is asserted.

    \item \textbf{Transformation Domain:}
          The range or class of transformations included
          in the claim.

    \item \textbf{Domain and Respect:}
          The analytical or ontological domain and the
          particular aspect in which preservation is
          being examined.

    \item \textbf{Preservation Criterion:}
          The criterion by which preservation is
          distinguished from alteration or other
          relations of comparison.

    \item \textbf{Admissibility Conditions:}
          The grounds for treating the transformation,
          comparison, and criterion as applicable to
          the claim.

    \item \textbf{Ontological Grounding:}
          The justification for attributing ontological
          significance to the determination and the
          preservation relation under examination.

    \item \textbf{Scope of Validity:}
          The conditions and limits within which the
          conclusion is warranted.
\end{enumerate}

These elements are analytical requirements for
evaluating invariance claims. They are not independent
ontological entities, nor do they impose one mandatory
formal notation upon every investigation.

Their purpose is to make the content of a claim
determinate and to prevent conclusions from extending
beyond what its specification and justification
support.

\subsection{The Forms of Invariance}

The chapter has examined several analytically
distinguishable forms of invariance:

\begin{itemize}
    \item Determinational Invariance;
    \item Structural Invariance;
    \item Relational Invariance;
    \item Compositional Invariance;
    \item Quantitative Invariance;
    \item Qualitative Invariance;
    \item Configurational Invariance;
    \item Relational-Structural Invariance.
\end{itemize}

These forms identify different respects in which
preservation may be examined. They do not constitute
a closed inventory of independent ontological
categories.

The forms may overlap. A transformation may preserve
a quantitative determination while also preserving
a structural or relational determination. Conversely,
a determination may be preserved in one respect
while being altered in another.

The classification is therefore analytical rather
than exhaustive. Its categories organize questions
about what is preserved; they do not establish that
every instance of preservation must belong to one
and only one form.

The identification of a particular form also does
not establish the presence of other forms. Structural
invariance does not automatically establish
compositional invariance, and the preservation of
a quantitative determination does not establish
the preservation of qualitative or relational
determinations.

The classification remains open to refinement
where further analysis identifies additional
respects of preservation or requires more precise
distinctions among existing forms.

\subsection{The Domains of Invariance}

The forms of invariance must be distinguished from
the domains within which preservation is examined.

The ontological framework developed in this book
distinguishes four proposed fundamental domains:

\begin{enumerate}
    \item First Determination;
    \item Structural Determination;
    \item Relational Manifestation;
    \item Material Manifestation.
\end{enumerate}

It also distinguishes three intermediate hubs:

\begin{enumerate}
    \item Hub of Decree;
    \item Hub of Signification;
    \item Hub of Realization.
\end{enumerate}

These four domains and three hubs are not seven
equivalent levels arranged in one undifferentiated
sequence. Their conceptual functions and relations
must be considered according to the broader
ontological framework. Their enumeration does
not, by itself, establish a linear order, temporal
sequence, causal hierarchy, or exhaustive taxonomy.

Within the analysis of invariance, the distinction
between forms and domains is essential. Forms concern
the respect of preservation being examined; domains
concern the ontological or analytical field within
which the relevant determination is considered.

An invariance claim associated with one domain
or hub cannot automatically be transferred to
another. Such a transfer requires an explicit
account of the determinations involved and of
the relation that makes the proposed correspondence
ontologically relevant.

The same restriction applies to formal and
ontological domains. A formal result may provide
a precise description of preservation within a
mathematical or logical system, but it does not
automatically establish that the corresponding
determination is ontologically instantiated.

\subsection{The Relations Between Invariance and Neighboring Concepts}

The chapter has established relations between
invariance and several neighboring concepts without
identifying them with one another.

\begin{itemize}
    \item \textbf{Identity:}
          Invariance concerns the preservation of a
          specified determination; identity concerns
          the existent as such. Neither the presence
          nor the absence of an invariant independently
          settles numerical identity.

    \item \textbf{Persistence:}
          Invariance concerns preservation under
          transformation; persistence concerns the
          continued existence of an existent across
          a specified range of conditions. An invariant
          does not independently establish persistence.

    \item \textbf{Continuity:}
          Invariance concerns a preserved determination;
          continuity concerns a relation connecting
          successive states or phases. Neither entails
          the other.

    \item \textbf{Difference and Distinction:}
          Invariance identifies a determination that
          remains preserved under specified conditions.
          Difference and distinction concern other
          relations among determinations or existents.
          Their criteria must not be collapsed into
          a single notion of change.

    \item \textbf{Structure:}
          Structural invariance concerns the preservation
          of a specified organizational determination.
          It does not establish the identity of the
          existent or the preservation of every component.

    \item \textbf{Composition:}
          Compositional invariance concerns the preservation
          of a specified compositional determination.
          It does not necessarily establish the persistence
          of every component or the numerical identity
          of the composed existent.

    \item \textbf{Stability:}
          Stability concerns the maintenance of a condition,
          configuration, or organization under specified
          circumstances or perturbations. It is not
          identical to invariance.

    \item \textbf{Immutability:}
          Immutability concerns the absence of change
          within a specified respect, whereas invariance
          concerns preservation under a specified
          transformation. An invariance claim does not
          imply unrestricted immutability.
\end{itemize}

These distinctions preserve the independent analytical
roles of the concepts required for examining the
existent and its transformations. They also prevent
the analysis from assigning explanatory priority
to invariance merely because preservation can be
identified in a particular case.

\subsection{The Limits of Invariance}

The limits established in the preceding section
govern the interpretation of every invariance claim.

No universal invariant is presumed to exist across
all existents and transformations. The identification
of an invariant in one case does not establish its
universal applicability.

The preservation of an invariant does not
independently establish numerical identity,
persistence, continuity, stability, immutability,
or the preservation of all other determinations.

The absence of an identified invariant does not
independently establish the absence or loss of
identity, the termination of persistence, or the
discontinuity of an existent or process.

Formal invariance does not automatically establish
ontological invariance. Observed preservation does
not justify conclusions extending beyond the
conditions under which it has been established.
Likewise, preservation in one domain does not
automatically establish preservation in another.

Finally, the inventory of identified invariants
cannot be treated as an exhaustive account of the
ontological constitution of an existent.

These restrictions follow from the defined scope
of invariance. They do not diminish its significance;
they prevent its analytical contribution from being
mistaken for a complete ontological explanation.

\subsection{Consolidated Framework of Invariance}

The preceding analysis can be consolidated through
a framework that distinguishes the elements of
an invariance claim from the forms and domains
within which preservation is investigated.

The framework must preserve three analytical
dimensions:

\begin{enumerate}
    \item \textbf{Claim Specification:}
          What is alleged to be preserved, under which
          transformation, and according to which
          criterion and conditions?

    \item \textbf{Form of Preservation:}
          Which determination or respect of preservation
          is being examined?

    \item \textbf{Domain of Examination:}
          Within which analytical or ontological domain
          is the claim situated?
\end{enumerate}

These dimensions are related but not interchangeable.
The form of invariance does not determine the domain,
and the domain does not, by itself, determine the
form of preservation. Neither dimension replaces
the need to specify the transformation and criterion.

% =============================================================================
% Project : THE_PHILOSOPHY_OF_VIABILITY
% Book    : Book I — Ontology of Viability
% Book ID : VS.B01
%
% File    : 4_ASSETS/tables/VS.B01.P01.C02.S11.SS10.P004.TL01.consolidated_framework_of_invariance.tex
%
% Asset   : TL01 — Consolidated Framework of Invariance
% Type    : Table
% Location: Section 11, Subsection 10, Paragraph 004
%
% Purpose :
%   To consolidate the three analytical dimensions of an
%   invariance claim: claim specification, form of preservation,
%   and domain of examination.
%
% Dependencies :
%   Standard LaTeX table environments
%
% =============================================================================

\begin{table}[htbp]
    \centering
    \small
    \renewcommand{\arraystretch}{1.4}
    \setlength{\tabcolsep}{5pt}

    \caption{Three analytical dimensions of an invariance claim}
    \label{tab:consolidated-framework-of-invariance}

    \begin{tabular}{
        |p{0.20\textwidth}
        |p{0.34\textwidth}
        |p{0.34\textwidth}|
    }
        \hline

        \textbf{Dimension}
        &
        \textbf{Principal Elements}
        &
        \textbf{Analytical Question}
        \\ \hline

        Claim Specification
        &
        Subject; specified determination; transformation;
        transformation domain; respect; criterion;
        admissibility conditions; ontological grounding;
        scope of validity
        &
        What is claimed to be preserved, under which
        transformation, and within what limits?
        \\ \hline

        Form of Preservation
        &
        Determinational; structural; relational;
        compositional; quantitative; qualitative;
        configurational; relational-structural
        &
        What kind of determination or respect is
        preserved?
        \\ \hline

        Domain of Examination
        &
        First Determination; Structural Determination;
        Relational Manifestation; Material Manifestation;
        Hub of Decree; Hub of Signification;
        Hub of Realization
        &
        Within which analytical or ontological domain
        is preservation examined?
        \\ \hline

    \end{tabular}
\end{table}

% =============================================================================
% END OF TABLE TL01
% =============================================================================

% =============================================================================
% Project : THE_PHILOSOPHY_OF_VIABILITY
% Book    : Book I — Ontology of Viability
% Book ID : VS.B01
%
% File    : 4_ASSETS/diagrams/VS.B01.P01.C02.S11.SS10.P004.DG01.dimensions_of_an_invariance_claim.tex
%
% Asset   : DG01 — Dimensions of an Invariance Claim
% Type    : Diagram
% Location: Section 11, Subsection 10, Paragraph 004
%
% Purpose :
%   To represent the distinction and coordination of claim
%   specification, form of preservation, and domain of examination.
%
% Requires (preamble):
%   \usepackage{tikz}
%   \usetikzlibrary{arrows.meta, positioning, fit}
% =============================================================================

\begin{figure}[htbp]
    \centering
    \hyphenpenalty=10000 \exhyphenpenalty=10000\relax  % no hyphenation inside diagram labels

    \begin{tikzpicture}[
        font=\small,
        node distance=0.4cm,
        % --- base box ---
        box/.style={
            draw, rounded corners=3pt, align=center,
            inner sep=6pt, line width=0.7pt
        },
        % --- node roles (names avoid clashes with built-in TikZ keys) ---
        dimbox/.style={box, text width=3.3cm, minimum height=4.1cm,
                       fill=blue!8, draw=blue!55!black},
        assessbox/.style={box, text width=5.4cm, minimum height=1.1cm,
                          font=\small\bfseries,
                          fill=blue!16, draw=blue!55!black},
        resultbox/.style={box, text width=5.4cm, minimum height=1.1cm,
                          dashed, fill=green!8, draw=green!40!black},
        % --- frames and labels ---
        vsgroup/.style={draw=black!45, dashed, rounded corners=5pt,
                        inner sep=8pt, line width=0.5pt},
        grouplabel/.style={font=\footnotesize\itshape, text=black!75,
                           align=center},
        % --- arrows ---
        directed/.style={-{Stealth[length=2.5mm]}, line width=0.9pt}
    ]

    % ---------------------------------------------------------------
    % Three analytical dimensions (equal rank, no order among them)
    % ---------------------------------------------------------------
    \node[dimbox] (specification) {
        \textbf{Claim Specification}\\[5pt]
        \footnotesize
        Subject, determination, transformation,
        respect, criterion, conditions, scope
    };
    \node[dimbox, right=of specification] (form) {
        \textbf{Form of Preservation}\\[5pt]
        \footnotesize
        Determinational, structural, relational, compositional,
        quantitative, qualitative, configurational and combined forms
    };
    \node[dimbox, right=of form] (domaindim) {
        \textbf{Domain of Examination}\\[5pt]
        \footnotesize
        Four proposed fundamental domains
        and three intermediate hubs
    };

    \node[vsgroup, fit=(specification)(form)(domaindim)] (dims) {};
    \node[grouplabel, above=0.1cm of dims] {
        Distinct analytical dimensions, jointly coordinated
    };

    % ---------------------------------------------------------------
    % Assessment and result
    % ---------------------------------------------------------------
    \node[assessbox, below=0.9cm of dims] (assessment) {
        Invariance Assessment
    };
    \node[resultbox, below=0.9cm of assessment] (result) {
        Preservation claim within a specified scope
    };

    \draw[directed] (dims)       -- (assessment);
    \draw[directed] (assessment) -- (result);

    % ---------------------------------------------------------------
    % Note
    % ---------------------------------------------------------------
    \node[below=0.7cm of result, text width=10cm, align=center,
          font=\footnotesize\itshape] {
        The assessment does not independently establish
        numerical identity, persistence, continuity,
        or the total ontological constitution of the existent.
    };

    \end{tikzpicture}

    \caption{The three analytical dimensions of an invariance
    claim. They coordinate the assessment of preservation
    without becoming interchangeable or establishing conclusions
    beyond the specified scope.}
    \label{fig:dimensions-of-an-invariance-claim}
\end{figure}

% =============================================================================
% END OF DIAGRAM DG01
% =============================================================================

The table and diagram summarize the analytical
architecture developed in this chapter. They are
not intended as an exhaustive taxonomy of all
possible invariance relations. Their function is
to preserve the distinctions required for a
well-specified claim and to prevent conclusions
from being transferred between analytical dimensions
without justification.

\subsection{Established Conclusions}

The following conclusions have been established
through the analysis of this chapter:

\begin{enumerate}
    \item Invariance is the preservation of a specified
          determination under a specified transformation,
          relative to a defined domain, respect, criterion,
          and scope of validity.

    \item Invariance is analytically distinct from
          ontological identity.

    \item The preservation of an invariant does not
          independently establish numerical identity.

    \item The alteration or absence of an identified
          invariant does not independently establish
          the loss or absence of identity.

    \item Invariance is distinct from persistence,
          continuity, stability, and immutability.

    \item Transformation may involve preservation
          and alteration in different determinations
          or respects.

    \item Structural and compositional invariance
          concern distinct analytical respects of
          preservation.

    \item Forms of invariance and domains of invariance
          are distinct analytical dimensions.

    \item Formal invariance does not automatically
          establish ontological invariance.

    \item An invariance claim must be restricted to
          the transformations, conditions, and criteria
          for which preservation has been established.

    \item No universal invariant is presumed across
          all existents and transformations.

    \item The identified invariants of an existent
          do not exhaust its ontological constitution.
\end{enumerate}

These conclusions establish the conceptual position
of invariance within the ontology of identity. They
do not imply that every question concerning the
ontological significance of preservation has been
resolved.

\subsection{Open Questions and Further Investigation}

Several questions remain open and require further
ontological investigation.

First, the precise relation between invariance and
the determinations constitutive of an existent
requires further clarification. The present chapter
has established that invariance cannot replace
identity, but it has not determined whether some
forms of preservation may have a more fundamental
ontological significance than others.

Second, the conditions under which invariance may
be established across multiple ontological domains
remain unresolved. The distinction between the
four fundamental domains and three intermediate
hubs has been maintained, but the relations that
could warrant cross-domain preservation claims
require further examination.

Third, the relation between invariance and extensive
transformation remains open. The chapter has
distinguished preservation from identity and has
rejected the reduction of identity to invariant
determinations. It has not, however, provided a
complete account of the ontological conditions
under which an existent may undergo extensive
transformation.

Fourth, the possibility of multiple identity-related
determinations within one existent requires further
investigation. The identification of multiple
invariants does not establish multiple identities,
just as the identification of one invariant does
not establish that an existent has only one identity.

Fifth, the relation between formal invariance and
ontological invariance requires more precise
criteria for determining when a formal preservation
corresponds to an ontological determination.

Sixth, the possible existence of invariants with
a broad or cross-contextual range of applicability
remains open. The rejection of a presumed universal
invariant does not exclude the possibility that
particular invariants may prove to have extensive
scope.

These questions are not conclusions that should
be resolved by stipulation. They identify the
boundaries of what the present analysis establishes
and indicate where further conceptual and ontological
work is required.

\subsection{The Ontological Significance of Invariance}

The contribution of invariance to the ontology
of identity lies in clarifying how preservation
may occur within transformation without being
mistaken for identity itself.

It establishes a way to distinguish determinations
that remain preserved from those that are altered
under specified conditions. This distinction makes
it possible to examine transformation without
treating every change as an undifferentiated
alteration of the existent in its entirety.

An existent may exhibit preservation and alteration
at the same time, provided that the determinations
and respects involved are distinguished. The
preservation of one determination does not cancel
the alteration of another; nor does the alteration
of one determination eliminate the possibility
that others remain preserved.

% =============================================================================
% Project : THE_PHILOSOPHY_OF_VIABILITY
% Book    : Book I — Ontology of Viability
% Book ID : VS.B01
%
% File    : 4_ASSETS/illustrations/VS.B01.P01.C02.S11.SS13.P002.IL01.preservation_and_alteration_within_transformation.tex
%
% Asset   : IL01 — Preservation and Alteration Within Transformation
% Type    : Illustration
% Location: Section 11, Subsection 13, Paragraph 002
%
% Purpose :
%   To illustrate the coexistence of preservation, alteration,
%   and undetermined status among different determinations
%   during the transformation of one existent.
%
% Requires (preamble):
%   \usepackage{tikz}
%   \usetikzlibrary{arrows.meta, calc, positioning, fit, backgrounds}
% =============================================================================

\begin{figure}[htbp]
    \centering

    \begin{tikzpicture}[
        font=\small,
        node distance=0.25cm,
        % --- one determination (base) ---
        item/.style={
            draw=black!70, rounded corners=3pt, align=center,
            text width=3.0cm, minimum height=0.7cm,
            inner sep=3pt, line width=0.8pt
        },
        % --- status styles (differ by line type AND fill) ---
        neutral/.style={item, fill=white},
        preserved/.style={item, fill=green!12, draw=green!40!black,
                          line width=1pt},
        altered/.style={item, fill=blue!10, draw=blue!55!black, dashed},
        undetermined/.style={item, fill=white, dotted, line width=1pt},
        % --- frames, titles, labels ---
        configbox/.style={draw=black!45, fill=gray!5, rounded corners=6pt,
                          inner sep=8pt, line width=0.7pt},
        cfgtitle/.style={font=\small\bfseries, align=center},
        vsgroup/.style={draw=black!45, dashed, rounded corners=5pt,
                        inner sep=8pt, line width=0.5pt},
        % --- arrow ---
        directed/.style={-{Stealth[length=2.5mm]}, line width=0.9pt}
    ]

    % ---------------------------------------------------------------
    % Initial configuration
    % ---------------------------------------------------------------
    \node[preserved] (a1) {Determination A};
    \node[neutral, below=of a1] (b1) {Determination B};
    \node[neutral, below=of b1] (c1) {Determination C};

    % ---------------------------------------------------------------
    % Resulting configuration (same row positions as the initial one)
    % ---------------------------------------------------------------
    \node[preserved, right=3.8cm of a1] (a2) {Determination A};
    \node[altered, below=of a2] (b2) {Determination B$'$};
    \node[undetermined, below=of b2] (c2) {Determination C?};

    % ---------------------------------------------------------------
    % Frames behind the items
    % ---------------------------------------------------------------
    \begin{scope}[on background layer]
        \node[configbox, fit=(a1)(b1)(c1)] (initial) {};
        \node[configbox, fit=(a2)(b2)(c2)] (result) {};
    \end{scope}

    \node[cfgtitle, above=0.15cm of initial] {Initial Configuration};
    \node[cfgtitle, above=0.15cm of result]  {Resulting Configuration};

    % ---------------------------------------------------------------
    % Transformation
    % ---------------------------------------------------------------
    \draw[directed] (initial.east) -- (result.west)
        node[midway, above=2pt, font=\small\bfseries] {Transformation};

    % ---------------------------------------------------------------
    % Key (centred under both configurations, same styles as the items)
    % ---------------------------------------------------------------
    \coordinate (mid) at ($(initial.south)!0.5!(result.south)$);

    \node[altered, text width=2.4cm, font=\footnotesize,
          below=0.9cm of mid] (kalt) {Altered};
    \node[preserved, text width=2.4cm, font=\footnotesize,
          left=0.3cm of kalt] (kpres) {Preserved};
    \node[undetermined, text width=2.4cm, font=\footnotesize,
          right=0.3cm of kalt] (kund) {Undetermined};

    \node[vsgroup, fit=(kpres)(kalt)(kund)] (key) {};

    % ---------------------------------------------------------------
    % Note
    % ---------------------------------------------------------------
    \node[below=0.6cm of key, text width=10cm, align=center,
          font=\footnotesize\itshape] {
        One determination remains preserved, another changes,
        and the status of a third has not been established.
        These different outcomes do not by themselves settle
        the numerical identity of the existent.
    };

    \end{tikzpicture}

    \caption{Preservation and alteration may coexist within
    one transformation. The diagram distinguishes a preserved
    determination, an altered determination, and one whose
    status remains undetermined. It makes no independent claim
    about the identity of the existent.}
    \label{fig:preservation-and-alteration-within-transformation}
\end{figure}

% =============================================================================
% END OF ILLUSTRATION IL01
% =============================================================================

This differentiated analysis is significant for
the examination of structure, composition, relations,
and material determinations. It allows the
investigation to identify what remains preserved,
what changes, and what has not yet been sufficiently
determined.

The significance of invariance therefore lies
neither in explaining the whole existent nor in
providing a universal criterion of identity. It
lies in making particular relations of preservation
within transformation conceptually explicit.

A properly delimited concept of invariance is
more precise and more useful than an unrestricted
one. Its value depends on the adequacy of the
determination, transformation, criterion, and
scope specified in each claim.

\subsection{Concluding Statement}

The central conclusion of this chapter is that
invariance is a determinate mode of preservation
within transformation, not the foundation or
exhaustive criterion of ontological identity.

A determination may remain preserved while other
determinations of the same existent undergo change.
An existent may undergo transformation without
every determination being altered. Conversely,
the preservation of a particular determination
across distinct existents does not establish
their numerical identity.

Neither preservation nor alteration, considered
in isolation, settles the ontological question
of identity.

The analysis of invariance must therefore maintain
a precise distinction between the determination
that is preserved, the transformation under which
it is preserved, and the existent whose identity
is under consideration.

By establishing this distinction, Chapter 02
defines the proper conceptual and inferential
scope of invariance within the ontology of identity.
It also identifies the questions that remain open
without converting them into unsupported principles.

The chapter thus provides a basis for further
examination of identity and manifestation while
preserving the distinction between the identity
of an existent and the particular determinations
that may remain invariant through its transformations.

% =============================================================================
% END OF SECTION 11 OF CHAPTER 02
% =============================================================================